\documentclass[12pt]{article}

\usepackage[utf8]{inputenc}
\usepackage[T1]{fontenc}
\usepackage{lmodern}
\usepackage[spanish]{babel}
\usepackage{booktabs}
\usepackage{amsmath}
\usepackage{forest}
\usepackage{float}
\usepackage{listings}
\usepackage{xcolor}
\usepackage{tikz}
\usetikzlibrary{arrows.meta, positioning, calc}


\definecolor{codegreen}{rgb}{0,0.6,0}
\definecolor{codegray}{rgb}{0.5,0.5,0.5}
\definecolor{codepurple}{rgb}{0.58,0,0.82}
\definecolor{backcolour}{rgb}{0.95,0.95,0.92}

% Definir colores para los nodos segun diferencia de valores
\definecolor{orangelight}{rgb}{1.0,0.8,0.6}
\definecolor{orangemedium}{rgb}{1.0,0.6,0.3}
\definecolor{orangedark}{rgb}{0.8,0.4,0.1}
\definecolor{bluelight}{rgb}{0.7,0.8,1.0}
\definecolor{bluemedium}{rgb}{0.4,0.6,1.0}
\definecolor{bluedark}{rgb}{0.2,0.4,0.8}

\lstdefinestyle{mystyle}{
    backgroundcolor=\color{backcolour},   
    commentstyle=\color{codegreen},
    keywordstyle=\color{magenta},
    numberstyle=\tiny\color{codegray},
    stringstyle=\color{codepurple},
    basicstyle=\ttfamily\footnotesize,
    breakatwhitespace=false,         
    breaklines=true,                 
    captionpos=b,                    
    keepspaces=true,                 
    numbers=left,                    
    numbersep=5pt,                  
    showspaces=false,                
    showstringspaces=false,
    showtabs=false,                  
    tabsize=2
}

\lstset{style=mystyle}

\sloppy
\setlength{\parindent}{0pt}

\begin{document}

% Título y materia
\begin{center}
  {\LARGE \textbf{Guía de Ejercicios\\Cross Validation}}\\[0.5em]
  {Analítica de Datos, Universidad de San Andrés}
\end{center}

Si encuentran algún error en el documento o hay alguna duda, mandenmé un mail a rodriguezf@udesa.edu.ar y lo revisamos.

\section{Ejercicios Conceptuales}

\subsection{Verdadero o Falso}

\begin{enumerate}
    \item Cross validation permite obtener una estimación más robusta del rendimiento del modelo.
    \item En k-fold cross validation, el dataset se divide en k particiones exactamente iguales.
    \item Leave-one-out cross validation es siempre la mejor opción para evaluar modelos.
    \item Cross validation se puede usar tanto para seleccionar hiperparámetros como para evaluar el rendimiento final.
    \item En 5-fold cross validation, cada observación se usa exactamente 4 veces para entrenar y 1 vez para validar.
    \item Stratified cross validation es importante para datasets desbalanceados.
    \item Cross validation siempre reduce el overfitting del modelo.
    \item Un mayor número de folds siempre es mejor en cross validation.
\end{enumerate}

\subsection{Opción Múltiple}

\begin{enumerate}
    \item ¿Cuál es la principal ventaja de cross validation sobre hold-out validation?
    \begin{enumerate}
        \item Es más rápido computacionalmente
        \item Usa toda la data para entrenamiento y validación
        \item No requiere dividir los datos
        \item Es más fácil de implementar
    \end{enumerate}
    
    \item En 10-fold cross validation, ¿qué porcentaje de los datos se usa para entrenar en cada fold?
    \begin{enumerate}
        \item 10\%
        \item 80\%
        \item 90\%
        \item 100\%
    \end{enumerate}
    
    \item ¿Cuándo es recomendable usar Leave-One-Out cross validation?
    \begin{enumerate}
        \item Con datasets muy grandes
        \item Con datasets muy pequeños
        \item Siempre que sea posible
        \item Nunca
    \end{enumerate}
    
    \item ¿Qué es stratified cross validation?
    \begin{enumerate}
        \item Una técnica para datasets grandes
        \item Una forma de mantener la proporción de clases en cada fold
        \item Un método más rápido de validación
        \item Una técnica solo para regresión
    \end{enumerate}
    
    \item ¿Cuál es el trade-off principal al elegir el número de folds?
    \begin{enumerate}
        \item Sesgo vs Varianza
        \item Velocidad vs Precisión
        \item Complejidad vs Simplicidad
        \item Todas las anteriores
    \end{enumerate}
\end{enumerate}

\section{Ejercicios Prácticos}

\subsection{Ejercicio 1 - Cálculo de Modelos a Entrenar}

Considerar los siguientes escenarios e indicar cuántos árboles es necesario entrenar en cada caso:

\begin{enumerate}
    \item En el entrenamiento de un modelo CART se usa 5-fold cross validation para seleccionar el hiperparámetro max\_depth de la siguiente grilla: [none, 10, 20, 30]
    
    \item En el entrenamiento de un modelo CART se usa 5-fold cross validation para seleccionar la mejor combinación de los siguientes hiperparámetros:
    \begin{itemize}
        \item max\_depth: [None, 10, 20, 30, 40, 50]
        \item min\_samples\_split: [2, 5, 10]
        \item min\_samples\_leaf: [1, 2, 4]
    \end{itemize}
    
    \item En el entrenamiento de un modelo Random Forest se usa 5-fold cross validation para seleccionar la mejor combinación de los siguientes hiperparámetros:
    \begin{itemize}
        \item n\_estimators: [100, 200, 500, 1000]
        \item max\_depth: [10, 20, 30]
        \item min\_samples\_split: [2, 5, 10]
    \end{itemize}
\end{enumerate}

\subsection{Ejercicio 2 - Interpretación de Resultados}

Tener en cuenta los siguientes resultados de la cross validation del hiperparámetro n\_estimators para un modelo CART. ¿Qué valor del hiperparámetro elegirían? ¿Por qué?

\begin{lstlisting}[language=Python]
{
'mean_fit_time': array([7.5903286, 12.5109745, 30.5498374, 61.4766215]),
'std_fit_time': array([2.8521561, 0.3776906, 0.5143247, 0.6323429]),
'mean_score_time': array([0.1880916, 0.0871274, 0.9984307, 1.9107995]),
'std_score_time': array([0.0091651, 0.0154946, 0.1213947, 0.1533748]),
'param_n_estimators': masked_array(data=[100, 200, 500, 1000]),
'params': [{'n_estimators': 100},
           {'n_estimators': 200},
           {'n_estimators': 500},
           {'n_estimators': 1000}],
'split0_test_score': array([0.7995629, 0.7997190, 0.7998751, 0.7998751]),
'split1_test_score': array([0.7995629, 0.7997190, 0.7997190, 0.7997190]),
'split2_test_score': array([0.7995316, 0.7998438, 0.7998438, 0.7998438]),
'split3_test_score': array([0.7993754, 0.7996877, 0.7996877, 0.7998438]),
'split4_test_score': array([0.7996877, 0.7998438, 0.7998438, 0.7998438]),
'mean_test_score': array([0.7995754, 0.799762, 0.7997939, 0.7998251]),
'std_test_score': array([1.2270604e-04, 6.7250609e-05, 7.5455567e-05, 5.4430172e-05]),
'rank_test_score': array([4, 3, 2, 1], dtype=int32)
}
\end{lstlisting}

\subsection{Ejercicio 3 - KNN y Cross Validation}

Considerar el siguiente conjunto de datos donde se quiere predecir la etiqueta, + o - en función de la variable X. Con este conjunto de datos se quiere entrenar un modelo knn, para el cual el hiperparámetro k se seleccionará por cross-validation:

\begin{table}[H]
    \centering
    \begin{tabular}{|c|c|c|c|c|c|c|c|c|c|c|}
        \hline
        & \textbf{1} & \textbf{2} & \textbf{3} & \textbf{4} & \textbf{5} & \textbf{6} & \textbf{7} & \textbf{8} & \textbf{9} & \textbf{10} \\
        \hline
        \textbf{X} & -0.1 & 0.7 & 1.0 & 1.6 & 2.0 & 2.5 & 3.0 & 3.5 & 4.1 & 4.9 \\
        \hline
        \textbf{Y} & - & + & + & - & + & + & - & - & + & + \\
        \hline
    \end{tabular}
\end{table}

\begin{enumerate}
    \item Dar el accuracy de cross-validation por leave-one-out (es decir 10-folds cross-validation) de 1 vecino más cercano para la muestra de entrenamiento.
    \item Dar el accuracy de cross-validation por leave-one-out (es decir 10-folds cross-validation) de 3 vecinos más cercanos para la muestra de entrenamiento.
    \item ¿Qué valor de k seleccionas?
\end{enumerate}

\subsection{Ejercicio 4 - Implementación de K-Fold}

Implementá k-fold cross validation desde cero en Python para evaluar un modelo de regresión lineal:

\begin{lstlisting}[language=Python]
import numpy as np
from sklearn.linear_model import LinearRegression
from sklearn.metrics import mean_squared_error

def k_fold_cv(X, y, k=5, model_class=LinearRegression):
    """
    Implementa k-fold cross validation
    
    Parameters:
    X: features
    y: target
    k: numero de folds
    model_class: clase del modelo a evaluar
    
    Returns:
    scores: lista de scores para cada fold
    """
    # Tu codigo aca
    pass

# Datos de ejemplo
X = np.random.randn(100, 3)
y = X[:, 0] + 2*X[:, 1] - X[:, 2] + np.random.randn(100)*0.1

# Evaluar modelo
scores = k_fold_cv(X, y, k=5)
print(f"MSE promedio: {np.mean(scores):.4f}")
print(f"Desviacion estandar: {np.std(scores):.4f}")
\end{lstlisting}

\subsection{Ejercicio 5 - Stratified Cross Validation}

Tenemos un dataset desbalanceado con las siguientes clases:

\begin{table}[H]
    \centering
    \begin{tabular}{cc}
        \toprule
        \textbf{Clase} & \textbf{Cantidad} \\
        \midrule
        A & 800 \\
        B & 150 \\
        C & 50 \\
        \bottomrule
    \end{tabular}
\end{table}

\begin{enumerate}
    \item ¿Por qué es importante usar stratified cross validation en este caso?
    \item ¿Cuántas observaciones de cada clase habría en cada fold si usamos 5-fold stratified CV?
    \item ¿Qué pasaría si usáramos k-fold simple sin estratificación?
\end{enumerate}

\subsection{Ejercicio 6 - Grid Search con Cross Validation}

Implementá grid search con cross validation para optimizar hiperparámetros de un SVM:

\begin{lstlisting}[language=Python]
from sklearn.svm import SVC
from sklearn.model_selection import GridSearchCV
from sklearn.datasets import make_classification

# Generar datos
X, y = make_classification(n_samples=1000, n_features=10, 
                          n_classes=2, random_state=42)

# Definir grilla de parametros
param_grid = {
    'C': [0.1, 1, 10, 100],
    'gamma': ['scale', 'auto', 0.001, 0.01, 0.1],
    'kernel': ['rbf', 'linear']
}

# Tu codigo aca para implementar grid search con CV
\end{lstlisting}

\subsection{Ejercicio 7 - Comparación de Modelos}

Usando cross validation, compará el rendimiento de los siguientes modelos en un problema de clasificación:

\begin{enumerate}
    \item Logistic Regression
    \item Random Forest
    \item SVM
    \item KNN
\end{enumerate}

¿Qué métrica usarías? ¿Cómo determinarías si las diferencias son estadísticamente significativas?

\subsection{Ejercicio 8 - Análisis de Varianza}

\begin{enumerate}
    \item ¿Por qué es importante reportar tanto la media como la desviación estándar de los scores de CV?
    \item ¿Qué indica una alta varianza en los scores de cross validation?
    \item ¿Cómo podríamos reducir la varianza en los resultados de CV?
\end{enumerate}

\newpage

\section{Anexo: Soluciones}

\subsection{Soluciones - Verdadero o Falso}

\begin{enumerate}
    \item \textbf{Verdadero.} CV proporciona múltiples estimaciones del rendimiento, reduciendo la dependencia de una única división train/test.
    \item \textbf{Falso.} Los folds pueden tener tamaños ligeramente diferentes si n no es divisible por k.
    \item \textbf{Falso.} LOOCV puede tener alta varianza y ser computacionalmente costoso.
    \item \textbf{Verdadero.} CV se usa tanto para model selection como para model evaluation.
    \item \textbf{Verdadero.} En 5-fold CV, cada observación se usa 4 veces para entrenar y 1 vez para validar.
    \item \textbf{Verdadero.} Mantiene la proporción de clases en cada fold.
    \item \textbf{Falso.} CV ayuda a detectar overfitting, pero no lo reduce directamente.
    \item \textbf{Falso.} Más folds aumenta la varianza y el costo computacional.
\end{enumerate}

\subsection{Soluciones - Opción Múltiple}

\begin{enumerate}
    \item \textbf{b)} Usa toda la data tanto para entrenamiento como para validación en diferentes momentos.
    \item \textbf{c)} 90\% (9 de los 10 folds se usan para entrenamiento).
    \item \textbf{b)} Con datasets muy pequeños, cuando queremos usar todos los datos posibles.
    \item \textbf{b)} Mantiene la proporción de clases en cada fold.
    \item \textbf{d)} Todas las anteriores: más folds aumentan sesgo pero reducen varianza, son más lentos pero más precisos, y más complejos pero potencialmente mejores.
\end{enumerate}

\subsection{Soluciones - Ejercicio 1}

\begin{enumerate}
    \item \textbf{CART con 5-fold CV y 4 valores de max\_depth:}
    \begin{itemize}
        \item 4 valores de hiperparámetro × 5 folds = \textbf{20 árboles}
    \end{itemize}
    
    \item \textbf{CART con 5-fold CV y múltiples hiperparámetros:}
    \begin{itemize}
        \item max\_depth: 6 valores
        \item min\_samples\_split: 3 valores  
        \item min\_samples\_leaf: 3 valores
        \item Total combinaciones: 6 × 3 × 3 = 54
        \item 54 combinaciones × 5 folds = \textbf{270 árboles}
    \end{itemize}
    
    \item \textbf{Random Forest con 5-fold CV:}
    \begin{itemize}
        \item n\_estimators: 4 valores
        \item max\_depth: 3 valores
        \item min\_samples\_split: 3 valores
        \item Total combinaciones: 4 × 3 × 3 = 36
        \item 36 combinaciones × 5 folds = 180 Random Forests
        \item Cada Random Forest tiene su propio n\_estimators (100, 200, 500, o 1000)
        \item Total de árboles individuales = 180 × promedio de n\_estimators
        \item Promedio de n\_estimators = (100+200+500+1000)/4 = 450
        \item \textbf{Total aproximado: 81,000 árboles individuales}
    \end{itemize}
\end{enumerate}

\subsection{Soluciones - Ejercicio 2}

Analizando los resultados:

\begin{itemize}
    \item \textbf{n\_estimators = 100:} score = 0.7996, tiempo = 7.59s, std = 1.23e-04
    \item \textbf{n\_estimators = 200:} score = 0.7998, tiempo = 12.51s, std = 6.73e-05
    \item \textbf{n\_estimators = 500:} score = 0.7998, tiempo = 30.55s, std = 7.55e-05
    \item \textbf{n\_estimators = 1000:} score = 0.7999, tiempo = 61.48s, std = 5.44e-05
\end{itemize}

\textbf{Recomendación: n\_estimators = 100}

\vspace{1em}

\textbf{Justificación:}
\begin{itemize}
    \item La diferencia de score entre 100 y 200 es solo 0.0002 (0.7996 vs 0.7998)
    \item Esta diferencia es muy pequeña y probablemente no sea estadísticamente significativa
    \item El tiempo aumenta 65\% (de 7.59s a 12.51s) para una ganancia mínima
    \item La desviación estándar es ligeramente mayor en 100, pero aún muy baja
    \item En términos prácticos, 0.0002 de diferencia no justifica el costo computacional adicional
\end{itemize}

\subsection{Soluciones - Ejercicio 3}

\textbf{Dataset ordenado por X:}
\begin{table}[H]
    \centering
    \begin{tabular}{|c|c|c|c|c|c|c|c|c|c|c|}
        \hline
        \textbf{Punto} & 1 & 2 & 3 & 4 & 5 & 6 & 7 & 8 & 9 & 10 \\
        \hline
        \textbf{X} & -0.1 & 0.7 & 1.0 & 1.6 & 2.0 & 2.5 & 3.0 & 3.5 & 4.1 & 4.9 \\
        \hline
        \textbf{Y} & - & + & + & - & + & + & - & - & + & + \\
        \hline
    \end{tabular}
\end{table}

\textbf{a) LOOCV con k=1 (1-NN):}

Para cada punto, encontramos su vecino más cercano y vemos si coincide la clase:

\begin{itemize}
    \item Punto 1 (X=-0.1, Y=-): vecino más cercano es punto 2 (Y=+) → \textbf{Error}
    \item Punto 2 (X=0.7, Y=+): vecino más cercano es punto 3 (Y=+) → \textbf{Correcto}  
    \item Punto 3 (X=1.0, Y=+): vecino más cercano es punto 2 (Y=+) → \textbf{Correcto}
    \item Punto 4 (X=1.6, Y=-): vecino más cercano es punto 5 (Y=+) → \textbf{Error}
    \item Punto 5 (X=2.0, Y=+): vecino más cercano es punto 4 (Y=-) → \textbf{Error}
    \item Punto 6 (X=2.5, Y=+): vecino más cercano es punto 5 (Y=+) → \textbf{Correcto}
    \item Punto 7 (X=3.0, Y=-): vecino más cercano es punto 6 (Y=+) → \textbf{Error}
    \item Punto 8 (X=3.5, Y=-): vecino más cercano es punto 7 (Y=-) → \textbf{Correcto}
    \item Punto 9 (X=4.1, Y=+): vecino más cercano es punto 8 (Y=-) → \textbf{Error}
    \item Punto 10 (X=4.9, Y=+): vecino más cercano es punto 9 (Y=+) → \textbf{Correcto}
\end{itemize}

\textbf{Accuracy k=1: 6/10 = 0.60}

\vspace{1em}

\textbf{b) LOOCV con k=3 (3-NN):}

Para cada punto, encontramos sus 3 vecinos más cercanos y usamos votación por mayoría:

\begin{itemize}
    \item Punto 1: vecinos [2,3,4] → clases [+,+,-] → mayoría + → \textbf{Error}
    \item Punto 2: vecinos [1,3,4] → clases [-,+,-] → mayoría - → \textbf{Error}
    \item Punto 3: vecinos [2,4,5] → clases [+,-,+] → mayoría + → \textbf{Correcto}
    \item Punto 4: vecinos [3,5,6] → clases [+,+,+] → mayoría + → \textbf{Error}
    \item Punto 5: vecinos [4,6,7] → clases [-,+,-] → mayoría - → \textbf{Error}
    \item Punto 6: vecinos [5,7,8] → clases [+,-,-] → mayoría - → \textbf{Error}
    \item Punto 7: vecinos [6,8,9] → clases [+,-,+] → mayoría + → \textbf{Error}
    \item Punto 8: vecinos [7,9,10] → clases [-,+,+] → mayoría + → \textbf{Error}
    \item Punto 9: vecinos [8,10,7] → clases [-,+,-] → mayoría - → \textbf{Error}
    \item Punto 10: vecinos [9,8,7] → clases [+,-,-] → mayoría - → \textbf{Error}
\end{itemize}

\textbf{Accuracy k=3: 1/10 = 0.10}

\vspace{1em}

\textbf{c) Selección de k:}
Elegiría \textbf{k=1} porque tiene mejor accuracy (0.60 vs 0.10). Con este dataset pequeño y ruidoso, k=3 incluye demasiados vecinos que añaden ruido a la predicción.

\subsection{Soluciones - Ejercicio 4}

\begin{lstlisting}[language=Python]
import numpy as np
from sklearn.linear_model import LinearRegression
from sklearn.metrics import mean_squared_error

def k_fold_cv(X, y, k=5, model_class=LinearRegression):
    """
    Implementa k-fold cross validation
    """
    n_samples = len(X)
    fold_size = n_samples // k
    indices = np.arange(n_samples)
    np.random.shuffle(indices)  # mezclar indices
    
    scores = []
    
    for i in range(k):
        # Definir indices de validacion
        start_idx = i * fold_size
        end_idx = start_idx + fold_size if i < k-1 else n_samples
        val_indices = indices[start_idx:end_idx]
        train_indices = np.concatenate([indices[:start_idx], 
                                       indices[end_idx:]])
        
        # Dividir datos
        X_train, X_val = X[train_indices], X[val_indices]
        y_train, y_val = y[train_indices], y[val_indices]
        
        # Entrenar modelo
        model = model_class()
        model.fit(X_train, y_train)
        
        # Predecir y evaluar
        y_pred = model.predict(X_val)
        score = mean_squared_error(y_val, y_pred)
        scores.append(score)
    
    return scores
\end{lstlisting}

\subsection{Soluciones - Ejercicio 5}

\begin{enumerate}
    \item \textbf{Importancia de stratified CV:} Con clases desbalanceadas, k-fold simple podría crear folds sin representación de clases minoritarias. Stratified CV garantiza que cada fold mantenga la proporción original de clases.
    
    \item \textbf{Observaciones por fold en 5-fold stratified CV:}
    \begin{itemize}
        \item Total: 1000 observaciones
        \item Proporción: A=80\%, B=15\%, C=5\%
        \item Por fold (200 obs): A=160, B=30, C=10
    \end{itemize}
    
    \item \textbf{Problema sin estratificación:} Algunos folds podrían tener muy pocas o ninguna observación de clases B y C, dando estimaciones sesgadas del rendimiento.
\end{enumerate}

\subsection{Soluciones - Ejercicio 6}

\begin{lstlisting}[language=Python]
from sklearn.svm import SVC
from sklearn.model_selection import GridSearchCV
from sklearn.datasets import make_classification

# Generar datos
X, y = make_classification(n_samples=1000, n_features=10, 
                          n_classes=2, random_state=42)

# Definir grilla de parametros
param_grid = {
    'C': [0.1, 1, 10, 100],
    'gamma': ['scale', 'auto', 0.001, 0.01, 0.1],
    'kernel': ['rbf', 'linear']
}

# Grid search con CV
grid_search = GridSearchCV(
    SVC(), 
    param_grid, 
    cv=5, 
    scoring='accuracy',
    n_jobs=-1,
    verbose=1
)

grid_search.fit(X, y)

print(f"Mejores parametros: {grid_search.best_params_}")
print(f"Mejor score: {grid_search.best_score_:.4f}")
\end{lstlisting}

\subsection{Soluciones - Ejercicio 7}

\begin{lstlisting}[language=Python]
from sklearn.model_selection import cross_val_score
from sklearn.linear_model import LogisticRegression
from sklearn.ensemble import RandomForestClassifier
from sklearn.svm import SVC
from sklearn.neighbors import KNeighborsClassifier
import numpy as np

# Modelos a comparar
models = {
    'Logistic Regression': LogisticRegression(),
    'Random Forest': RandomForestClassifier(),
    'SVM': SVC(),
    'KNN': KNeighborsClassifier()
}

# Comparar modelos
results = {}
for name, model in models.items():
    scores = cross_val_score(model, X, y, cv=5, scoring='f1_weighted')
    results[name] = scores
    print(f"{name}: {scores.mean():.4f} (+/- {scores.std()*2:.4f})")

# Test estadistico (t-test pareado) para comparar modelos
from scipy import stats
model1_scores = results['Random Forest']
model2_scores = results['SVM'] 
t_stat, p_value = stats.ttest_rel(model1_scores, model2_scores)
print(f"p-value: {p_value:.4f}")
\end{lstlisting}

Usaría \textbf{F1-score weighted} para datasets desbalanceados o \textbf{accuracy} para balanceados. Para significancia estadística, usaría t-test pareado entre scores de CV.

\subsection{Soluciones - Ejercicio 8}

\begin{enumerate}
    \item \textbf{Importancia de media y desviación:} La media indica el rendimiento esperado, la desviación indica la estabilidad del modelo. Ambas son cruciales para tomar decisiones informadas.
    
    \item \textbf{Alta varianza indica:}
    \begin{itemize}
        \item Modelo sensible a cambios en datos de entrenamiento
        \item Posible overfitting
        \item Dataset pequeño o muy heterogéneo
        \item Necesidad de regularización
    \end{itemize}
    
    \item \textbf{Reducir varianza:}
    \begin{itemize}
        \item Aumentar el dataset
        \item Usar modelos más simples
        \item Aplicar regularización
        \item Aumentar el número de folds (trade-off con costo computacional)
        \item Usar stratified CV para datasets desbalanceados
    \end{itemize}
\end{enumerate}

\end{document}
